% !TEX root = ../main.tex

\sysusetup{
  keywords  = {异常检测, 候选召回, FlashANNS, GPU 并行计算, 金融风控},
  keywords* = {Outlier Detection, Candidate Retrieval, FlashANNS, GPU Parallel Computing, Financial Risk Control},
}

\begin{abstract}
金融大规模异常检测既要识别低异常比例数据中的高风险对象，也要满足近实时业务对吞吐和尾延迟的要求。样本规模扩大后，全量距离计算和异常评分会产生高代价中间结果；索引超出显存时，SSD访问进入关键路径；CPU--GPU往返及多设备共享PCIe、NUMA和I/O路径还会放大阶段边界开销。本文研究如何在给定检测方法与单机GPU环境下压缩这些成本。

本文构建FlashTOD“候选召回---异常评分”两阶段链路，这是全文的核心系统创新。FlashANNS从全量样本中召回候选集合$C_i$，其规模$|C_i|$受候选预算约束；PyTOD在候选范围内进行张量化精细评分，将评分输入由全量样本规模$n$缩减为$|C_i|$。两阶段还需转换数据语义：系统依据FlashANNS返回的候选编号、近邻距离和局部排序，恢复原始样本与候选集合的索引映射，并构造PyTOD所需的候选张量、距离张量和评分中间张量。接口连接只是这段执行链路的一部分。

单GPU实现处理这条链路的数据通路和运行时开销。查询向量、候选结果与评分张量尽量驻留设备侧，异步I/O与计算重叠、缓冲区复用和批次执行用于减少CPU--GPU往返、重复内存申请和阶段同步。多GPU实现承担并行与规模扩展：replicated架构使各GPU独立完成本地召回和评分，CPU负责请求分发和小规模结果汇总；容量或路径受限时，系统再采用设备侧归并、通信和拓扑感知I/O控制。请求分配与负载均衡属于调度，数据划分、设备计算和结果归并分属并行执行与跨设备结果处理。

实验以约1M规模的真实或半真实金融数据验证有效性，以10M、50M和100M高保真合成数据测试吞吐、延迟与扩展趋势。在本文的数据规模、硬件平台和候选预算下，FlashTOD的PR-AUC为0.906，PyTOD全量评分为0.932；QPS由2700提高到4280，增幅约58.5\%。单卡完整优化将QPS从3180提高到4550，p99延迟由67.0 ms降至45.5 ms。replicated路径下，1、2、4 GPU的QPS分别为4550、7780和13620，4 GPU吞吐约为单GPU的2.99倍。

候选预算带来检测效果与执行效率的折中。FlashTOD改进执行链路和并行实现，不改变PyTOD评分模型；结论限于本文数据集、硬件平台和单机多GPU环境。
\end{abstract}

\begin{abstract*}
Large-scale financial outlier detection must identify rare high-risk objects while meeting near-real-time throughput and tail-latency requirements. At larger sample scales, full-distance computation and anomaly scoring create expensive intermediate results; SSD access enters the critical path when an index exceeds GPU memory; CPU--GPU round trips and contention on shared PCIe, NUMA, and I/O paths further amplify stage-boundary costs. This thesis studies how these costs can be reduced for a fixed scoring method on a single GPU-equipped server.

This thesis presents FlashTOD as its central system contribution: a two-stage anomaly detection pipeline for large-scale financial data. FlashANNS performs ANNS-based candidate retrieval and returns a candidate set $C_i$ whose size $|C_i|$ is bounded by the candidate budget. PyTOD then performs fine-grained tensorized scoring over these candidates, reducing its input scope from the full sample size $n$ to $|C_i|$. A direct API connection is insufficient: candidate IDs, neighbor distances, local ranks, and reference/index mappings must be converted into the candidate tensors, distance tensors, and scoring intermediates consumed by PyTOD. This device-side data contract translates retrieval results into anomaly-scoring semantics.

Single-GPU data-path optimization keeps query vectors, candidate results, and scoring tensors on the device whenever possible. Asynchronous I/O--compute overlap, buffer reuse, and batched execution reduce CPU--GPU transfers, repeated allocation, and stage synchronization. Multi-GPU parallel execution uses a replicated architecture in which each GPU completes local retrieval and scoring, while the CPU distributes requests and aggregates small outputs. Capacity- or path-constrained cases additionally use device-side merging, communication, and topology-aware I/O control. Request assignment and load balancing are scheduling functions; data partitioning, device computation, and merging belong to parallel execution and cross-device result processing.

Experiments use real or semi-real financial data at around 1M scale for effectiveness evaluation and high-fidelity synthetic data at 10M, 50M, and 100M scales for throughput, latency, and scalability tests. Under the evaluated data scales, hardware platform, and candidate budgets, FlashTOD obtains a PR-AUC of 0.906, compared with 0.932 for PyTOD full scoring, while QPS increases from 2700 to 4280 (\,+58.5\%). The full single-GPU path raises QPS from 3180 to 4550 and reduces p99 latency from 67.0 ms to 45.5 ms. In the replicated path, 1, 2, and 4 GPUs deliver 4550, 7780, and 13620 QPS, respectively; four GPUs reach about 2.99$\times$ the single-GPU throughput.

The results expose a trade-off between detection quality and execution efficiency controlled by the candidate budget. FlashTOD changes the execution pipeline and its parallel implementation rather than the PyTOD scoring model. The conclusions are limited to the evaluated datasets, hardware platform, and single-node multi-GPU environment.
\end{abstract*}
